\section{Extension triadique d’Euler et conservation des formes}
\label{eul:capitulo}

La représentation quadratique de TRIT acquiert un contenu opératoire en
identifiant les générateurs sélectionnés par APP et par le calendrier du
TPK. Ce chapitre reconstruit ces antécédents avant d’exprimer leurs
réalisations exponentielles. L’ordre distingue l’opération discrète,
sa représentation matricielle et la fonction analytique qui permet de décrire
son évolution en un paramètre réel.

Trois constructions interviennent : le cycle multiplicatif d’ordre trois
dans APP, le transport élémentaire qui conserve une coordonnée de
prolongation et l’incidence entre modes qui découle du calendrier
tritique. Leurs générateurs satisfont respectivement les relations
elliptique, parabolique et hyperbolique. La forme alternée du plan permet
ensuite d’établir une conservation commune, en conservant expressément les différences métriques
et d’orientation de chaque réalisation.

\subsection{Cycle multiplicatif d’APP et générateur elliptique}
\label{eul:coxeter}

Dans \(\mathbb Z/9\mathbb Z\), la multiplication par quatre est inversible
et d’ordre trois :
\[
 4^2=7\pmod9,\qquad 4^3=1\pmod9.
\]
Sur les unités apparaissent les cycles \((1,4,7)\) et \((2,8,5)\).
Les classes \(0,3,6\) restent fixes. La sélection du multiplicateur et de ses
orbites appartient donc à l’arithmétique nonadique d’APP.

\subsubsection{Action contragrédiente sur les deux évaluations}

Soient \(\bar S(x,y)=x+y\) et \(\bar P(x,y)=xy\), à valeurs dans
\(\mathbb Z/9\mathbb Z\). Pour \(a=4\),
\begin{equation}
 \bar S(ax,ay)=a\bar S(x,y),\qquad
 \bar P(ax,ay)=a^2\bar P(x,y)=a^{-1}\bar P(x,y).
\label{eul:eq:contragrediente}
\end{equation}
La première identité découle de la distributivité ; la seconde, de
\(a^3=1\). La somme et le produit parcourent des orientations opposées du
même cycle résiduel. Dans les évaluations entières, les quotients
nonadiques restent associés au relèvement de ces classes.

\subsubsection{Module d’augmentation et matrice du cycle}

Un des cycles \((r,4r,7r)\) étant fixé, considérons le module libre
engendré par \(e_r,e_{4r},e_{7r}\). Son sous-module d’augmentation est
\[
 L=\{x_0e_r+x_1e_{4r}+x_2e_{7r}:x_0+x_1+x_2=0\}.
\]
Une base est \(b_1=e_r-e_{4r}\), \(b_2=e_{4r}-e_{7r}\). Le cycle
envoie \(b_1\) sur \(b_2\) et \(b_2\) sur \(-b_1-b_2\). Sa matrice,
agissant sur les vecteurs colonnes, est
\[
 C=\begin{pmatrix}0&-1\\1&-1\end{pmatrix}.
\]
Le produit scalaire hérité du module à trois coordonnées a pour matrice
de Gram
\[
 G_{\mathrm e}=\begin{pmatrix}2&-1\\-1&2\end{pmatrix}.
\]
On obtient ainsi le réseau \(A_2\), avec sa forme déterminée par les
différences de positions du cycle.

\begin{proposicion}[Normalisation elliptique du cycle APP]
\label{eul:prop:eliptico}
Le cycle satisfait
\[
 C^2+C+I=0,\qquad C^\mathsf TG_{\mathrm e}C=G_{\mathrm e}.
\]
Dans \(L\otimes_{\mathbb Z}\mathbb R\), l’opérateur
\[
 J_{+1}=\frac{2C+I}{\sqrt3}
\]
est de trace nulle et satisfait \(J_{+1}^2=-I\).
\end{proposicion}
\begin{proof}
La multiplication donne
\(C^2=\bigl(\begin{smallmatrix}-1&1\\-1&0\end{smallmatrix}\bigr)\),
d’où \(C^2+C+I=0\). La permutation cyclique conserve le
produit scalaire des trois coordonnées et son sous-module d’augmentation ;
dans la base choisie, cela équivaut à l’identité de Gram. Enfin,
\((2C+I)^2=4C^2+4C+I=-3I\), et
\(\operatorname{tr}(2C+I)=0\).
\end{proof}

La normalisation réelle fournit une racine opératorielle de \(-I\) à
partir du cycle déjà construit. L’identité peut également s’écrire
\[
 C=-\frac12I+\frac{\sqrt3}{2}J_{+1}.
\]
Son interprétation angulaire est obtenue après avoir fixé la réalisation
analytique du type elliptique.

\subsection{Transport de seuil et générateur parabolique}
\label{eul:parabolico}

Considérons les coordonnées entières \((x,c)\), où \(c\) est conservé
comme incrément enregistré. Le transport élémentaire
\[
 (x,c)\longmapsto(x+c,c)
\]
est représenté par
\[
 U_0=I+N,
 \qquad N=\begin{pmatrix}0&1\\0&0\end{pmatrix}.
\]
C’est le modèle linéaire du transport de seuil utilisé dans la
réalisation quadratique. La division équilibrée et sa normalisation,
développées au chapitre~\ref{trit:capitulo}, conservent leurs règles
propres lorsqu’elles sont appliquées à une évaluation entière.

\begin{proposicion}[Itération et réversibilité du transport]
\label{eul:prop:parabolico}
On a \(N^2=0\), \(N\ne0\) et, pour tout \(m\in\mathbb Z\),
\[
 U_0^m=I+mN,
 \qquad U_0^m(x,c)=(x+mc,c).
\]
Le transport conserve le déterminant et admet l’inverse \(I-N\).
\end{proposicion}
\begin{proof}
L’image de \(N\) est contenue dans la première coordonnée, sur laquelle
\(N\) s’annule ; d’où \(N^2=0\). L’identité
\((I+aN)(I+bN)=I+(a+b)N\) prouve les puissances positives, l’inverse et
les puissances négatives. Le déterminant de \(I+mN\) est un.
\end{proof}

L’indice parabolique \(\tau=0\) correspond à un générateur de carré
nul. L’opérateur \(N\) transporte une donnée linéaire et peut être différent
de zéro sur un état concret. La signature nulle, l’opérateur nul et un
générateur nilpotent sont trois objets distincts.

\subsection{Incidence des modes dérivée du calendrier TPK}
\label{eul:incidencia-modos}

La construction de la section~\ref{sec:cal-incidencia} fournit
la matrice~\eqref{eq:cal-incidencia}. Son action primitive est explicitée ici
pour relier le décompte tritique à la représentation quadratique
hyperbolique développée ci-dessous.

Le sélecteur tritique de phase prend, sur les positions \(1,\ldots,9\),
les valeurs \((+1,-1,0)\) répétées trois fois. Sur le mode opératoire,
on utilise la règle
\[
 +1:\ m\mapsto\Sigma,\qquad
 -1:\ m\mapsto\Pi,\qquad
 0:\ m\mapsto m.
\]
Le bloc primitif produit, après chaque mise à jour, la succession
de modes \((\Sigma,\Pi,\Pi)\). Le transport spatial et la mise à jour
du registre demeurent dans la transition complète du TPK ; on lit ici
son incidence modale.

\begin{teorema}[Incidence primitive du calendrier]
\label{eul:teo:fav}
En comptant les paires consécutives avec clôture cyclique dans
\((\Sigma,\Pi,\Pi)\), la matrice obtenue dans l’ordre \((\Sigma,\Pi)\)
est
\[
 F_{\rm av}=\begin{pmatrix}0&1\\1&1\end{pmatrix}.
\]
Dans les calendriers de longueur \(9,27,108\), les décomptes sont
\(3F_{\rm av},9F_{\rm av},36F_{\rm av}\), respectivement.
\end{teorema}
\begin{proof}
Les trois paires sont \(\Sigma\to\Pi\), \(\Pi\to\Pi\) et
\(\Pi\to\Sigma\). Chacune apparaît une fois. La paire
\(\Sigma\to\Sigma\) apparaît zéro fois. Cela détermine les quatre
entrées. La répétition du calendrier multiplie tous les décomptes
par le nombre de blocs primitifs, qui est \(3,9,36\) dans les cas
indiqués.
\end{proof}

Le lecteur de feuillets associe \(\Sigma\) à la lecture surfacique et \(\Pi\) à
la lecture volumique. Cette attribution transporte l’incidence obtenue vers
les deux types de lecture. La matrice résulte du calendrier ; son analyse
spectrale s’effectue une fois ses entrées construites.

Les décomptes de calendriers répétés et les puissances matricielles ont
des fonctions différentes. \(36F_{\rm av}\) compte les paires parcourues
pendant un calendrier de \(108\) événements. En revanche,
\(F_{\rm av}^{108}\) compte les chaînes admissibles de longueur \(108\)
dans le graphe d’incidence. Le premier objet suit une chronologie ;
le second appartient à son enveloppe d’adjacences.

\subsubsection{Réalisation hyperbolique de l’incidence}

Deux compositions d’incidence produisent
\[
 A=F_{\rm av}^2=\begin{pmatrix}1&1\\1&2\end{pmatrix}.
\]
Ses coefficients comptent les parcours de deux arêtes. On a
\(\det A=1\), \(\operatorname{tr}A=3\) et, par multiplication directe,
\(A^2-3A+I=0\).

\begin{proposicion}[Générateur hyperbolique et forme invariante]
\label{eul:prop:hiperbolico}
L'opérateur
\[
 J_{-1}=\frac{2A-3I}{\sqrt5}
\]
est de trace nulle et de carré \(I\). La forme symétrique
\[
 G_{\mathrm h}=\begin{pmatrix}2&1\\1&-2\end{pmatrix}
\]
a pour signature \((1,1)\) et satisfait
\(A^\mathsf TG_{\mathrm h}A=G_{\mathrm h}\).
\end{proposicion}
\begin{proof}
De \(A^2-3A+I=0\), on obtient
\((2A-3I)^2=4A^2-12A+9I=5I\). La trace du numérateur est nulle.
De plus,
\[
 G_{\mathrm h}A=\begin{pmatrix}3&4\\-1&-3\end{pmatrix},
 \qquad
 A^\mathsf TG_{\mathrm h}A
 =\begin{pmatrix}2&1\\1&-2\end{pmatrix}.
\]
Le déterminant de \(G_{\mathrm h}\) est \(-5\), donc ses deux
valeurs propres sont de signes opposés.
\end{proof}

La réalisation est lorentzienne en dimension \(1+1\) : sont spécifiés
un plan réel, une forme de signature \((1,1)\) et un transport qui la
conserve. La réalisation physique d’un espace-temps utilise ensuite
ses champs, dimensions et applications de transduction correspondants.

\subsection{Identité exponentielle des trois régimes}
\label{eul:identidad}

Les trois opérateurs concrets agissent dans des espaces de provenance distincte :
le module d’augmentation d’un cycle APP, le plan du transport de seuil
et le plan d’incidence modale. Leurs normalisations satisfont
\[
 J_\tau^2=-\tau I,\qquad \operatorname{tr}J_\tau=0,
 \qquad \tau\in\{+1,0,-1\}.
\]
Cette relation détermine leurs puissances et permet de construire les fonctions
d’évolution par des séries convergentes.

\begin{definicion}[Fonctions du régime quadratique]
Pour \(t\in\mathbb R\), on définit
\[
 c_\tau(t)=\sum_{k=0}^{\infty}
       \frac{(-\tau)^kt^{2k}}{(2k)!},\qquad
 s_\tau(t)=\sum_{k=0}^{\infty}
       \frac{(-\tau)^kt^{2k+1}}{(2k+1)!}.
\]
Ces fonctions réalisent analytiquement le type quadratique déjà obtenu.
\end{definicion}

\begin{teorema}[Extension triadique de l’identité d’Euler]
\label{eul:teo:euler}
Pour les générateurs précédents,
\begin{equation}
 E_\tau(t):=\exp(tJ_\tau)
 =c_\tau(t)I+s_\tau(t)J_\tau.
\label{eul:eq:euler}
\end{equation}
Les spécialisations de \((c_\tau,s_\tau)\) sont
\((\cos,\sin)\), \((1,t)\) et \((\cosh,\sinh)\), dans l’ordre
\(\tau=+1,0,-1\). De plus,
\[
 c_\tau(t)^2+\tau s_\tau(t)^2=1,
 \qquad E_\tau(t)^{-1}=E_\tau(-t).
\]
\end{teorema}
\begin{proof}
La série exponentielle converge absolument en norme matricielle.
La relation quadratique donne
\(J_\tau^{2k}=(-\tau)^kI\) et
\(J_\tau^{2k+1}=(-\tau)^kJ_\tau\). Séparer les termes pairs et
impairs produit \eqref{eul:eq:euler}. Les trois spécialisations se
reconnaissent dans leurs séries. Comme les matrices \(tJ_\tau\) et \(-tJ_\tau\)
commutent, leurs exponentielles ont pour produit \(I\). La parité des
séries donne
\[
 E_\tau(t)E_\tau(-t)
 =\bigl(c_\tau(t)^2+\tau s_\tau(t)^2\bigr)I,
\]
ce qui prouve l’identité scalaire et l’inverse.
\end{proof}

Les nombres qui apparaissent lors de l’évaluation de fonctions analytiques ont ici
le statut de reconnaissance d’opérateurs préalablement construits.
La génération régionale coinductive conserve son émetteur, ses relèvements
et ses cylindres. En particulier, ces séries décrivent la réalisation
exponentielle ; la sélection des préfixes régionaux s’effectue au moyen
du TPK.

\subsubsection{Loi de composition et invariants quadratiques}

La commutation des multiples d’un même générateur donne
\(E_\tau(t)E_\tau(s)=E_\tau(t+s)\). En comparant les coefficients de
\(I,J_\tau\),
\begin{align}
 c_\tau(t+s)&=c_\tau(t)c_\tau(s)-\tau s_\tau(t)s_\tau(s),\\
 s_\tau(t+s)&=s_\tau(t)c_\tau(s)+c_\tau(t)s_\tau(s).
\end{align}
La conjugaison quadratique change \(t\) en \(-t\) au sein du régime
fixé. L’inversion orientée d’une succession tritique, quant à elle,
peut changer les indices de régime et doit être transportée par
l’involution correspondante du TPK.

Dans le type elliptique, l’identité conservée est \(a^2+b^2=1\).
Dans le type hyperbolique, \(a^2-b^2=1\), et les coordonnées \(u=a+b\),
\(v=a-b\) vérifient \(uv=1\). L’évolution hyperbolique a
\((u,v)=(e^t,e^{-t})\) : l’expansion d’une coordonnée s’accompagne
de la contraction de l’autre. Dans le type parabolique,
\(E_0(t)=I+tN\), avec un transport linéaire et un déterminant égal à un.

\subsubsection{Évaluations sur les opérateurs discrets}

La représentation du cycle elliptique donne
\[
 C=E_{+1}(2\pi/3),\qquad
 C^n=E_{+1}(2\pi n/3),\quad n\in\mathbb Z.
\]
En particulier,
\[
 E_{+1}(\pi)+I=0,\qquad E_{+1}(2\pi)=I,
 \qquad E_0(1)=I+N.
\]
La coordonnée angulaire \(\pi\) exprime analytiquement le demi-tour ;
son identification à la sortie régionale HMT relève du lecteur
coinductif développé à sa place causale.

L’incidence \(F_{\rm av}\) a pour polynôme caractéristique
\(x^2-x-1\). Sa racine positive \(\lambda\) satisfait
\(\lambda^2=\lambda+1\) ; en la reconnaissant comme nombre d’or, on écrit
\(\lambda=\varphi\). Avec \(\eta=2\log\lambda\),
\[
 \cosh\eta=\frac32,\qquad \sinh\eta=\frac{\sqrt5}{2},
 \qquad E_{-1}(\eta)=A=F_{\rm av}^2.
\]
La vérification utilise \(\lambda^{-1}=\lambda-1\), d’où
\(\lambda^2+\lambda^{-2}=3\) et
\(\lambda^2-\lambda^{-2}=\sqrt5\). La matrice a été obtenue auparavant par
décompte des transitions ; la reconnaissance spectrale caractérise son
échelle. L’égalité avec la coordonnée régionale engendrée est conservée
avec le lecteur et la démonstration de cette région.

De même, l’identité scalaire \(\exp(I)=eI\) caractérise l’évaluation
unitaire de l’exponentielle. Sa fonction dans la représentation analytique
est distincte de la construction coinductive de la région correspondante.

\subsection{Résolution polynomiale des régimes}
\label{eul:proyectores}

Dans la somme directe des trois espaces réels, définissons
\[
 \mathbb J=J_{+1}\oplus N\oplus J_{-1},\qquad Q=\mathbb J^2.
\]
Sur ces termes, \(Q\) agit comme \(-I,0,I\), respectivement.
Donc \(Q^3=Q\) et \(\mathbb J^6=\mathbb J^2\).

\begin{proposicion}[Idempotents de régime]
\label{eul:prop:proyectores}
Les opérateurs
\[
 p_{\mathrm e}=\frac{Q^2-Q}{2},\qquad
 p_{\mathrm p}=I-Q^2,\qquad
 p_{\mathrm h}=\frac{Q^2+Q}{2}
\]
sont des idempotents orthogonaux au sens algébrique et leur somme vaut \(I\).
Leurs images sont les trois termes de la représentation.
\end{proposicion}
\begin{proof}
Les trois polynômes prennent les valeurs \((1,0,0)\), \((0,1,0)\) et
\((0,0,1)\) sur \(-1,0,+1\). Leurs carrés moins eux-mêmes et
leurs produits croisés s’annulent en ces trois points, ils sont donc
multiples de \(X^3-X\). La substitution de \(Q\) démontre l’idempotence et
l’annulation des produits croisés. La somme des trois polynômes
est un. Leur action sur chaque bloc identifie les images.
\end{proof}

L’exponentielle complète s’exprime par
\begin{align}
 \exp(t\mathbb J)={}&p_{\mathrm e}(\cos t\,I+\sin t\,\mathbb J)
 +p_{\mathrm p}(I+t\mathbb J)\notag\\
 &+p_{\mathrm h}(\cosh t\,I+\sinh t\,\mathbb J).
\end{align}
Cette résolution distingue exactement les régimes. Les transitions
entre eux relèvent des applications du TPK sur les états avec orientation,
mode et mémoire ; la somme directe conserve leurs domaines séparés.

\subsection{Compression nonadique et registre complémentaire}
\label{eul:compresion}

La conservation entre une lecture et son registre complémentaire découle
de la structure à neuf phases. Soit \(V\) un espace muni d’un produit
scalaire et \(C:V\to V\) un transport linéaire. Dans \(V^9\), définissons
\[
 S_C(v_0,\ldots,v_8)=(Cv_8,v_0,\ldots,v_7),
 \qquad Jv=\frac13(v,\ldots,v).
\]
L’inclusion \(J\) est isométrique parce qu’elle réunit neuf copies avec le facteur
\(1/3\). Son adjoint somme les composantes et divise par trois. La
compression sur cette inclusion est
\[
 T=J^*S_CJ=\frac{8I+C}{9}.
\]
Huit composantes conservent \(v\) et une traverse le transport \(C\).
Le complément orthogonal de l’inclusion enregistre
\[
 (I-JJ^*)S_CJv
 =\frac1{27}\bigl(8(C-I)v,-(C-I)v,\ldots,-(C-I)v\bigr).
\]
Le carré de sa norme est
\(8\|(C-I)v\|^2/81\). Le même contenu vectoriel peut être représenté
isométriquement par
\[
 D=\frac{\sqrt8}{9}(C-I),
\]
en maintenant le registre à neuf composantes lorsque l’on utilise leurs
étiquettes individuelles.

\begin{proposicion}[Récupération algébrique à partir de deux lectures]
\label{eul:prop:recuperacion}
Pour tout \(v\in V\),
\begin{equation}
 v=Tv-\frac1{\sqrt8}Dv.
\label{eul:eq:recuperacion}
\end{equation}
Si \(C\) est isométrique, on a en outre
\[
 \|v\|^2=\|Tv\|^2+\|Dv\|^2.
\]
\end{proposicion}
\begin{proof}
La première identité résulte de
\((8I+C-(C-I))/9=I\). Pour la seconde, en développant les produits,
on obtient
\[
 T^*T+D^*D=\frac{72I+9C^*C}{81}.
\]
L'hypothèse \(C^*C=I\) donne l'identité. Les termes mixtes s'annulent avec les coefficients \(+8\) et \(-8\).
\end{proof}

La récupération utilise les registres vectoriels, y compris leurs corrélations. Les intensités \(\|Tv\|^2\) et \(\|Dv\|^2\) fournissent le bilan scalaire, tandis que la formule inverse utilise \(Tv\) et \(Dv\) dans leur intégralité.

\subsection{Conservation de l'aire orientée}
\label{eul:area}

Dans tout plan réel orienté, soit
\[
 \Omega=\begin{pmatrix}0&1\\-1&0\end{pmatrix}.
\]
Pour une matrice \(L=\bigl(\begin{smallmatrix}a&b\\c&d\end{smallmatrix}\bigr)\), la multiplication directe donne
\begin{equation}
 L^\mathsf T\Omega L=(ad-bc)\Omega=(\det L)\Omega.
\label{eul:eq:determinante-area}
\end{equation}
Les générateurs tritiques étant de trace nulle, \(\det E_\tau(t)=1\). Leurs réalisations préservent donc la forme alternée, bien que leurs formes métriques diffèrent.

\begin{teorema}[Conservation nonadique commune aux trois régimes]
\label{eul:teo:area} Pour \(C=E_\tau(t)\), soient \(T=(8I+C)/9\) et \(D=\sqrt8(C-I)/9\). Alors
\begin{equation}
 \Omega=T^\mathsf T\Omega T+D^\mathsf T\Omega D.
\label{eul:eq:balance-area}
\end{equation}
\end{teorema}
\begin{proof}
Le développement donne
\[
 T^\mathsf T\Omega T+D^\mathsf T\Omega D
 =\frac{72\Omega+9C^\mathsf T\Omega C}{81}.
\]
Les termes \(C^\mathsf T\Omega\) et \(\Omega C\) s'annulent. Par \eqref{eul:eq:determinante-area}, \(C^\mathsf T\Omega C=\Omega\), et le résultat est \(\Omega\).
\end{proof}

Avec \(c=c_\tau(t)\), les contributions déterminantielles sont
\begin{equation}
 \det T=\frac{65+16c}{81},\qquad
 \det D=\frac{16(1-c)}{81},\qquad
 \det T+\det D=1.
\label{eul:eq:determinantes}
\end{equation}
En effet, pour \(\det C=1\), \(\det(aI+C)=a^2+a\operatorname{tr}C+1\), et \(\operatorname{tr}C=2c\).

\begin{tablacompleta}
\caption{Contributions à l'aire orientée de trois transports engendrés.}
\begin{tabular}{lccc}
\toprule
Transporte&\(\operatorname{tr}C\)&\(\det T\)&\(\det D\)\\
\midrule Cycle APP \(C\)&$-1$&$19/27$&$8/27$\\
Seuil \(I+N\)&$2$&$1$&$0$\\
Incidence \(F_{\rm av}^2\)&$3$&$89/81$&$-8/81$\\
\bottomrule
\end{tabular}
\notatabla{Chaque paire a pour somme un. Le signe enregistre l'orientation ; le tableau concerne l'aire alternée, et non les probabilités.}
\end{tablacompleta}

Dans le transport parabolique, \(D\) peut être non nul et de rang un : son aire est nulle et son registre contient une information linéaire. Dans le transport hyperbolique, le complément a l'orientation opposée. Dans le transport elliptique, la forme positive propre au cycle permet en outre le bilan de norme. Les trois réalisations partagent l'identité alternée en conservant leurs types.

\subsection{Formes variables et restitution entre plans successifs}
\label{eul:formas-variables}

Considérons une suite de plans \(V_n\) munis de matrices symétriques définies positives \(G_n\). Soient \(B_n,C_n:V_n\to V_{n+1}\) deux transports satisfaisant
\[
 B_n^\mathsf TG_{n+1}B_n=G_n,
 \qquad C_n^\mathsf TG_{n+1}C_n=G_n.
\]
Le transport \(B_n\) identifie les coordonnées entre les deux plans ; \(C_n\) inclut l'évolution supplémentaire. Définissons
\[
 T_n=\frac{8B_n+C_n}{9},\qquad
 D_n=\frac{\sqrt8}{9}(C_n-B_n).
\]

\begin{teorema}[Conservation et récupération avec une forme variable]
\label{eul:teo:forma-variable}
On a
\begin{align}
 G_n&=T_n^\mathsf TG_{n+1}T_n+D_n^\mathsf TG_{n+1}D_n,
 \label{eul:eq:balance-variable}\\
 v&=B_n^{-1}\left(T_nv-\frac1{\sqrt8}D_nv\right).
 \label{eul:eq:recuperacion-variable}
\end{align}
\end{teorema}
\begin{proof}
Dans le développement de la première expression, les termes croisés s'annulent. Il reste \(72B_n^\mathsf TG_{n+1}B_n/81\) et \(9C_n^\mathsf TG_{n+1}C_n/81\), dont la somme vaut \(G_n\). L'identité \(T_n-D_n/\sqrt8=B_n\) prouve la seconde formule. La positivité des formes rend \(B_n\) injectif ; entre plans de même dimension, il est inversible.
\end{proof}

Pour \(I_n(v)=\tfrac12v^\mathsf TG_nv\), \(v_{n+1}=T_nv_n\) et \(m_n=D_nv_n\), on obtient
\[
 I_n(v_n)=I_{n+1}(v_{n+1})+I_{n+1}(m_n).
\]
La somme sur \(N\) étapes donne
\begin{equation}
 I_0(v_0)=I_N(v_N)+\sum_{n=0}^{N-1}I_{n+1}(m_n).
\label{eul:eq:telescopica}
\end{equation}
Chaque registre est évalué dans la forme du plan où il a été produit. Les transports inverses de \eqref{eul:eq:recuperacion-variable} restituent tous les états dans l'ordre inverse à partir de l'état terminal et des registres complets.

Une réalisation concrète utilise des matrices inversibles \(M_n\), \(G_n=M_n^{-\mathsf T}M_n^{-1}\) et des rotations \(R_n\) :
\[
 B_n=M_{n+1}M_n^{-1},\qquad
 C_n=M_{n+1}R_nM_n^{-1}.
\]
Les deux identités de formes se vérifient en substituant ces expressions et en utilisant \(R_n^\mathsf TR_n=I\). Les coordonnées, les échelles et les paramètres de \(M_n,R_n\) proviennent des lecteurs de l'état lorsque cette réalisation est appliquée à la Mécanique Dimensionnelle.

\subsubsection{Critère exact de transfert au registre}

Si \(R_n\) est une rotation d'angle \(\theta_n\), dans les coordonnées normalisées \(y_n=M_n^{-1}v_n\) on a
\[
 y_{n+1}=\frac{8I+R_n}{9}y_n,
 \qquad
 \left(\frac{8I+R_n}{9}\right)^\mathsf T
 \left(\frac{8I+R_n}{9}\right)=r_nI,
\]
\[
 r_n=1-\frac{32}{81}\sin^2(\theta_n/2).
\]
Par récurrence, \(I_N(v_N)=I_0(v_0)\prod_{n<N}r_n\).

\begin{proposicion}[Critère d'annulation de l'état terminal]
Si \(I_0(v_0)>0\), l'état terminal satisfait \(I_N(v_N)\to0\) exactement lorsque
\[
 \sum_{n\ge0}\sin^2(\theta_n/2)=+\infty.
\]
\end{proposicion}
\begin{proof}
Pour \(c=32/81\) et \(u\in[0,1]\), l'intégration de \(1/(1-s)\) entre zéro et \(cu\) donne
\[
 cu\le-\log(1-cu)\le\frac{cu}{1-c}.
\]
La somme des logarithmes diverge exactement lorsque la somme des \(u_n=\sin^2(\theta_n/2)\) diverge. Le produit des \(r_n\) tend vers zéro exactement dans ce cas.
\end{proof}

La conservation \eqref{eul:eq:telescopica} vaut à toute profondeur finie. Le critère précédent détermine quand la totalité de la forme initiale est transférée au registre dans cette réalisation. La suite angulaire est celle obtenue par le lecteur correspondant de la trajectoire.

\subsection{Action intégrée et compatibilité avec le raffinement}
\label{eul:accion-refinamiento}

Dans un plan canonique de coordonnées \((Q,P)\), soit \(\gamma\) une courbe fermée \(C^1\) par morceaux. Pour toute matrice réelle \(L\) d'ordre deux,
\begin{equation}
 \oint_{L\gamma}P\,\mathrm dQ
 =\det(L)\oint_\gamma P\,\mathrm dQ.
\label{eul:eq:accion-determinante}
\end{equation}
Pour le démontrer, écrivons \(Q'=aQ+bP\), \(P'=cQ+dP\). Les intégrales fermées de \(Q\,\mathrm dQ\) et \(P\,\mathrm dP\) sont nulles. L'intégration de \(\mathrm d(PQ)\) donne \(\oint Q\,\mathrm dP=-\oint P\,\mathrm dQ\). Il reste le coefficient \(ad-bc\), qui prouve la formule.

Appliquée à \eqref{eul:eq:determinantes}, elle donne
\[
 \oint_\gamma P\,\mathrm dQ
 =\oint_{T\gamma}P\,\mathrm dQ
  +\oint_{D\gamma}P\,\mathrm dQ.
\]
L'orientation de chaque courbe préserve le signe de sa contribution.

Le raffinement exige en outre la compatibilité entre les opérateurs de niveaux successifs. Dans cette section, soient \(\iota_n:V_n\to V_{n+1}\) des injections linéaires et \(A_n:V_n\to V_n\) des endomorphismes satisfaisant \(A_{n+1}\iota_n=\iota_nA_n\). Nous définissons à chaque niveau \(T_n=(8I_{V_n}+A_n)/9\) et \(D_n=\sqrt8(A_n-I_{V_n})/9\). Ces expressions impliquent
\[
 T_{n+1}\iota_n=\iota_nT_n,
 \qquad D_{n+1}\iota_n=\iota_nD_n.
\]
La restitution \eqref{eul:eq:recuperacion} commute alors avec le raffinement. Cette identité conserve simultanément l'opération et le registre ; sa vérification dans une réalisation concrète utilise l'application \(\iota_n\) produite par ce prolongement.

\subsubsection{Conservation du facteur scalaire du lecteur}

Si une réalisation reçoit un facteur d'échelle \(e^{-x}\), le transport est \(C=e^{-x}E_\tau(t)\). Son déterminant vaut \(e^{-2x}\), et le même développement de la loi nonadique donne
\[
 T^\mathsf T\Omega T+D^\mathsf T\Omega D
 =\frac{8+e^{-2x}}9\Omega.
\]
Pour \(x\ge0\), le complément d'échelle
\[
 D_{\rm esc}=\frac{\sqrt{1-e^{-2x}}}{3}I
\]
restitue le bilan complet :
\[
 \Omega=T^\mathsf T\Omega T+D^\mathsf T\Omega D
       +D_{\rm esc}^\mathsf T\Omega D_{\rm esc}.
\]
Le coefficient s'obtient à partir du défaut déterminantiel de l'échelle. Chaque application conserve le facteur de son lecteur et précise si ce registre supplémentaire est nécessaire.

L'extension triadique d'Euler s'articule ainsi avec des opérations engendrées, des formes invariantes et des opérateurs de récupération. Le cycle APP, le transport nilpotent et l'incidence modale conservent leur provenance ; la réalisation analytique décrit leurs évolutions ; la division nonadique de la lecture conserve le complément nécessaire à la restitution. Cette succession exprime le lien entre régime, transformation et mémoire sans remplacer la généalogie par une formule exponentielle isolée.
